\section{Estruturas Sedimentares}

\subsection{Definição}

As \textbf{estruturas sedimentares} são feições desenvolvidas durante ou após a
deposição dos sedimentos e que podem permanecer preservadas no depósito ou na
rocha resultante de sua litificação \parencite{harms1975,nichols2009}.

Essas estruturas constituem registros dos processos responsáveis pela formação
do depósito e permitem interpretar aspectos como:

\begin{itemize}
    \item direção e sentido das paleocorrentes;
    \item regime e energia do fluxo;
    \item mecanismos de transporte e deposição;
    \item exposição subaérea;
    \item processos químicos ocorridos durante ou após a deposição;
    \item presença e atividade de organismos;
    \item deformação de sedimentos ainda inconsolidados.
\end{itemize}

De acordo com sua origem, podem ser agrupadas em:

\begin{enumerate}
    \item estruturas físicas ou mecânicas;
    \item estruturas químicas;
    \item estruturas biogênicas.
\end{enumerate}

\subsection{Estruturas físicas ou mecânicas}

As \textbf{estruturas físicas} resultam principalmente da ação mecânica dos
agentes de transporte, deposição ou deformação dos sedimentos, como correntes,
ondas, marés, vento e gravidade.

\subsubsection{Estratificação}

A \textbf{estratificação} corresponde à organização dos sedimentos em camadas ou
estratos. Cada estrato pode representar um episódio de deposição ou um intervalo
durante o qual determinadas condições sedimentares permaneceram relativamente
constantes.

As superfícies que separam os estratos são denominadas
\textbf{planos de estratificação}. Mudanças entre camadas podem refletir
variações de granulometria, composição, cor, energia do ambiente e velocidade de
sedimentação.

\subsubsection{Laminação}

A \textbf{laminação} corresponde à organização do sedimento em camadas muito
delgadas. As lâminas podem resultar de pequenas variações de granulometria,
composição ou condições de fluxo.

Quando apresentam geometria aproximadamente horizontal e paralela, constituem a
\textbf{laminação plano-paralela}.

\subsubsection{Estratificação cruzada}

A \textbf{estratificação cruzada} ocorre quando conjuntos de lâminas encontram-se
inclinados em relação ao plano principal de estratificação.

Sua formação está normalmente associada à migração de formas de leito, como
marcas onduladas, dunas subaquosas, dunas eólicas e barras. A orientação das
lâminas pode fornecer informações sobre a direção do paleofluxo.

Em ambientes com reversão periódica de corrente, conjuntos inclinados em sentidos opostos podem formar \textbf{estratificação cruzada espinha de peixe}. Superfícies de reativação registram erosão parcial e retomada da migração de uma forma de leito, frequentemente após mudança na intensidade ou na direção do fluxo.

\subsubsection{Marcas onduladas}

As \textbf{marcas onduladas}, ou \textit{ripple marks}, são pequenas ondulações
formadas na superfície de sedimentos inconsolidados pela ação da água ou do
vento.

As marcas produzidas por corrente são normalmente assimétricas e podem indicar o
sentido predominante do fluxo. As produzidas pela ação oscilatória das ondas
tendem a apresentar maior simetria.

\subsubsection{Acamamentos flaser, ondulado e lenticular}

Quando areia e lama se alternam sob variações periódicas da energia, podem formar-se três padrões principais:

\begin{itemize}
  \item \textbf{flaser}: lâminas descontínuas de lama preservadas entre ripples predominantemente arenosos;
  \item \textbf{ondulado} (\textit{wavy}): proporções semelhantes de areia e lama formam camadas lateralmente contínuas e onduladas;
  \item \textbf{lenticular}: lentes isoladas de areia ficam envolvidas por lama.
\end{itemize}

Esses padrões são comuns em depósitos submetidos à alternância entre correntes capazes de transportar areia e períodos de baixa energia que permitem a decantação de lama \parencite{nichols2009}.

\subsubsection{Estratificação cruzada hummocky e swaley}

A estratificação cruzada \textit{hummocky} apresenta laminações suavemente curvas, com elevações e depressões de baixo ângulo. É associada principalmente a fluxos combinados produzidos por ondas e correntes durante tempestades em ambientes marinhos rasos. A variante \textit{swaley} é dominada por depressões erosivas e tende a ocorrer em condições de maior retrabalhamento. A interpretação deve considerar a associação de fácies e não apenas uma geometria isolada \parencite{harms1975,nichols2009}.

\subsubsection{Estratificação gradacional}

A \textbf{estratificação gradacional} ocorre quando há variação sistemática do
tamanho dos grãos ao longo de uma camada.

Na gradação normal:

\[
\text{base grossa}
\rightarrow
\text{topo fino}.
\]

Esse padrão geralmente indica diminuição progressiva da energia do fluxo durante
a deposição. Também pode ocorrer gradação inversa, na qual o tamanho dos grãos
aumenta em direção ao topo.

\subsubsection{Marcas de sola}

As \textbf{marcas de sola} são feições preservadas principalmente na base de
camadas sedimentares. Podem ser formadas quando correntes erodem um substrato
ainda inconsolidado e as depressões resultantes são posteriormente preenchidas
por sedimento.

Algumas dessas estruturas permitem interpretar direção e sentido de
paleocorrentes.

\subsubsection{Gretas de dessecação}

As \textbf{gretas de dessecação} formam-se quando sedimentos finos, úmidos e
coesivos são expostos ao ar e perdem água. A retração do material produz uma rede
de fraturas frequentemente poligonais.

Sua presença indica exposição subaérea temporária e pode ocorrer em planícies de
maré, margens de lagos, planícies de inundação e outras áreas sujeitas a
alternância entre inundação e secagem.

\subsubsection{Estruturas de carga e deformações sinsedimentares}

As \textbf{estruturas de carga} podem desenvolver-se quando uma camada mais densa
é depositada sobre outra menos densa e ainda pouco consolidada. A camada
superior pode afundar parcialmente no sedimento inferior.

Sedimentos saturados em água também podem sofrer deformações antes da
litificação em consequência de escorregamentos, sobrecarga, liquefação, escape
de fluidos ou instabilidade gravitacional.

Podem resultar:

\begin{itemize}
    \item dobras convolutas;
    \item laminações deformadas;
    \item estruturas de escape de fluidos;
    \item escorregamentos sinsedimentares.
\end{itemize}

\subsection{Estruturas químicas}

As \textbf{estruturas químicas} resultam principalmente de processos de
dissolução, precipitação, recristalização ou substituição mineral.

Podem formar-se durante a deposição ou durante a diagênese e são controladas por
fatores como:

\begin{itemize}
    \item concentração de íons;
    \item salinidade;
    \item pH;
    \item condições de oxidação e redução;
    \item circulação de fluidos;
    \item evaporação.
\end{itemize}

\subsubsection{Cristais e moldes de evaporitos}

Em ambientes sujeitos à intensa evaporação, a concentração de sais dissolvidos
pode alcançar a saturação e provocar a precipitação de minerais como gipsita,
anidrita e halita.

Os cristais podem posteriormente ser dissolvidos, deixando moldes que podem ser
preenchidos por outros minerais ou por sedimento. Quando a forma original é
preservada após a substituição mineral, forma-se um \textbf{pseudomorfo}.

\subsubsection{Concreções}

As \textbf{concreções} são corpos formados pela precipitação localizada de
minerais nos espaços porosos de um sedimento ou rocha sedimentar.

Podem crescer ao redor de um núcleo, como um fóssil, fragmento mineral, matéria
orgânica ou irregularidade do sedimento. Podem ser compostas, por exemplo, por
calcita, sílica, óxidos de ferro ou sulfetos.

\subsubsection{Nódulos}

Os \textbf{nódulos} são concentrações minerais localizadas dentro de uma rocha
sedimentar e podem apresentar composição diferente da matriz que os envolve.
Exemplos incluem nódulos de sílex, fosfato e materiais carbonáticos.

\subsubsection{Estruturas de dissolução}

A circulação de fluidos pode promover dissolução seletiva de minerais,
originando cavidades, superfícies irregulares e porosidade secundária. Essas
cavidades podem ser posteriormente preenchidas por outros minerais.

\subsubsection{Estilólitos}

Os \textbf{estilólitos} são superfícies irregulares, geralmente serrilhadas,
formadas principalmente por dissolução sob pressão.

São particularmente comuns em rochas carbonáticas e estão associados a
processos diagenéticos relacionados ao soterramento.

\subsection{Estruturas biogênicas}

As \textbf{estruturas biogênicas} são produzidas direta ou indiretamente pela
atividade dos organismos.

Podem registrar atividades de deslocamento, alimentação, escavação, habitação,
repouso, perfuração e construção. Muitas dessas estruturas são estudadas pela
\textbf{icnologia}.

\subsubsection{Bioturbação}

A \textbf{bioturbação} corresponde ao retrabalhamento do sedimento pela atividade
dos organismos.

Os organismos podem escavar galerias, revolver o substrato, ingerir partículas,
deslocar materiais e misturar diferentes camadas. Uma bioturbação intensa pode
destruir parcial ou completamente estruturas sedimentares anteriormente
existentes.

\subsubsection{Icnofósseis}

Os \textbf{icnofósseis} são registros preservados da atividade de organismos, em
vez de partes do próprio corpo.

Entre os exemplos encontram-se:

\begin{itemize}
    \item pegadas;
    \item pistas;
    \item rastros;
    \item galerias;
    \item escavações;
    \item marcas de repouso;
    \item perfurações.
\end{itemize}

Os icnofósseis possuem grande importância paleoambiental porque registram a
relação entre os organismos e o substrato.

\subsubsection{Escavações, galerias e perfurações}

Organismos que vivem no interior do sedimento podem produzir galerias
horizontais, verticais, inclinadas ou ramificadas.

As \textbf{perfurações} diferem das escavações porque são produzidas em
substratos relativamente duros, como rochas, conchas, fragmentos carbonáticos ou
superfícies consolidadas.

\subsubsection{Pegadas, pistas e rastros}

O deslocamento de organismos sobre superfícies sedimentares pode produzir marcas
posteriormente soterradas e preservadas.

Essas estruturas podem fornecer informações sobre a direção do movimento, o
comportamento dos organismos e as características do substrato.

\subsubsection{Estromatólitos}

Os \textbf{estromatólitos} são estruturas organossedimentares laminadas
associadas principalmente à atividade de comunidades microbianas.

Tapetes microbianos podem aprisionar partículas, ligar sedimentos e favorecer a
precipitação mineral. A repetição desses processos produz estruturas laminadas
de crescimento progressivo.

\subsubsection{Recifes e construções biogênicas}

Alguns organismos participam diretamente da construção de corpos sedimentares.
Corais, algas calcárias, esponjas e outros organismos podem formar recifes e
outras construções carbonáticas.

Essas estruturas podem controlar localmente a circulação da água, a deposição de
sedimentos, a distribuição dos organismos e o desenvolvimento de diferentes
fácies sedimentares.

\subsection{Estruturas primárias e secundárias}

As estruturas sedimentares também podem ser classificadas conforme o momento de
sua formação.

As \textbf{estruturas primárias} desenvolvem-se durante ou imediatamente após a
deposição. Entre os exemplos estão:

\begin{itemize}
    \item estratificação;
    \item laminação;
    \item estratificação cruzada;
    \item marcas onduladas;
    \item gradação granulométrica;
    \item marcas de sola.
\end{itemize}

As \textbf{estruturas secundárias} formam-se posteriormente, geralmente durante
deformações precoces ou processos diagenéticos. Entre elas estão:

\begin{itemize}
    \item estruturas de carga;
    \item estruturas de escape de fluidos;
    \item concreções;
    \item nódulos;
    \item estilólitos;
    \item estruturas de dissolução e substituição.
\end{itemize}

\subsection{Importância para a interpretação paleoambiental}

As estruturas sedimentares constituem registros dos processos que atuavam no
ambiente durante a formação dos depósitos.

Por exemplo:

\begin{itemize}
    \item estratificação cruzada pode indicar direção de paleocorrente;
    \item marcas onduladas assimétricas podem indicar o sentido do fluxo;
    \item marcas onduladas simétricas podem indicar ação de ondas;
    \item gretas de dessecação indicam exposição subaérea;
    \item gradação normal pode indicar redução progressiva da energia de um fluxo;
    \item evaporitos e seus moldes podem indicar forte evaporação;
    \item bioturbação registra atividade de organismos no substrato;
    \item estromatólitos registram interação entre atividade microbiana e sedimentação.
\end{itemize}

A interpretação de uma estrutura isolada deve ser realizada com cautela. A
reconstrução do ambiente deposicional torna-se mais confiável quando são
analisadas conjuntamente a granulometria, a composição mineralógica, as
estruturas sedimentares, os fósseis, a geometria dos corpos sedimentares e as
relações estratigráficas.

De maneira geral:

\[
\boxed{
\begin{gathered}
\text{intemperismo}
\rightarrow
\text{erosão}
\rightarrow
\text{transporte}
\\
\rightarrow
\text{deposição}
\rightarrow
\text{estruturas sedimentares}
\rightarrow
\text{diagênese}
\end{gathered}
}
\]


% =========================================================
% DIAGÊNESE
% =========================================================
